% ============================================================
%  ch7_conclusion_future.tex
%  Chapter 7: Conclusions and Future Work
%  Master's Thesis -- SQL MATCH_RECOGNIZE on Pandas DataFrames
% ============================================================
% !TEX root = Thesis.tex
\chapter{Conclusions and Future Work}
\label{chap:conclusions}

This chapter brings the findings together.  It interprets the results,
states the limits of the work, and identifies the next steps.
Section~\ref{sec:discussion} discusses the evidence.
Section~\ref{sec:conc-validation} maps the main claims to that evidence,
and Section~\ref{sec:conc-limitations} defines their boundaries.
Sections~\ref{sec:conclusion} and~\ref{sec:conc-future} give the final
conclusion and future work.

% ----------------------------------------------------------------
\section{Discussion}
\label{sec:discussion}
% ----------------------------------------------------------------

The results show that a practical subset of SQL row-pattern recognition
can run directly on \texttt{pandas} DataFrames.  For the reported
correctness queries, the engine agreed with the applicable Trino~473,
Oracle, and hand-written \texttt{pandas} baselines.
The cases covered partitioning, ordering, implicit \texttt{DEFINE}
conditions, navigation, match numbering, and measure output.  This
supports the evaluated subset.  It does not claim full SQL:2016 or
vendor parity.

The usability comparison gives a second result.  The three SQL examples
used 14--16 lines of detection logic, compared with 24--32 lines for
their hand-written \texttt{pandas} versions.  Those functions also needed
explicit grouping, shifting, candidate filtering, index tracking, and
skip handling.  The SQL form keeps this logic in one declarative
statement.  That lowers developer effort, makes the pattern easier to
read, and travels between notebook analysis and SQL workflows---as long
as the query stays inside the supported subset.

The performance benchmark contained 30 pattern--size
combinations: five pattern families over six sizes from 100\,K to the
full 2.22\,M rows.  Each cell is the mean of twenty measured runs after
five warmups and the stated outlier rule.  Each system used the same
input and query logic under one CPU core and the same outer 32\,GB
memory limit.  The
proposed engine averaged
6\,110\,392~rows/s and 0.17\,s per cell.  Oracle averaged 1.04\,s and
Trino 1.98\,s.  The proposed engine had the lowest time in all 30
cells, and all 30 normalized outputs agreed across the three systems.
Pattern structure remained important:
\texttt{complex\_nested} was slower than the simpler families.

Memory shows the trade-off.  Across the main benchmark, the maximum
incremental query peaks were 189.05\,MB for the engine, 58.71\,MB for
Oracle, and 20.75\,MB for Trino.  The maximum total footprints showed
the reverse order: 523.97\,MB, 2{,}057.17\,MB, and
14{,}992.90\,MB, respectively.  All six values use the same cgroup-based
metric, but the process boundaries differ: one engine process versus a
complete database container.  Vendor-native counters have different
scopes and are not treated as equivalent measurements.

This result has a clear boundary.  The engine is specialized for one
in-process DataFrame workload.  It avoids database planning, storage,
and client-transfer layers, and it compiles supported conditions to
column operations.  The evaluated running-average predicate is also
updated incrementally.  Supported expressions outside a compiled form
use the general evaluator; unsupported forms are rejected.
State-dependent conditions and some ambiguous pattern
choices may require exact search.  In the separate state-dependent
test, the engine processed the full 2.22\,M rows in 1.28\,s and
preserved the reference output.  These results describe the tested
subset on one machine.  They are not a general ranking of SQL systems
(Sections~\ref{subsec:eval-findings} and
\ref{subsec:eval-worst-cases}).

A separate scalability test used Amazon Reviews~2023 from 5\,M to
227.9\,M rows.  Execution time was close to linear for all three
systems.  The engine's measured memory was also close to linear.  On
the size averages, the engine remained
$5.1$--$10.1\times$ faster than Oracle~21c EE and
$11.4$--$13.8\times$ faster than Trino over the four pattern families
completed by all three systems.  Oracle~21c EE completed every
cell.  Trino's \texttt{complex\_nested} pattern reached the 600\,s
limit at the last two sizes.  The engine completed the full corpus, but
its maximum process footprint reached about 35.7\,GB.  Speed and memory
residency must therefore be read together.  The engine values describe
the fixed benchmark revision recorded in Appendix~\ref{app:provenance}.
The retained Oracle and Trino measurements came from separate campaigns
under the same stated protocol.  They were not run simultaneously.

The architecture separates parsing, validation, logical planning,
automata construction, matching, and output assembly.  The optimized
paths use row-local masks, compact DFA and run-length plans, compiled
measures, fused output, bounded caches, lazy row access, and conservative
output projection.  Partition processing remains sequential.  The
study therefore evaluates a single-node \texttt{pandas} engine, not a
distributed SQL system.  Trino and Oracle retain wider SQL coverage,
mature optimizers, and production execution infrastructure.

% ----------------------------------------------------------------
\section{Validation of Chapter Claims}
\label{sec:conc-validation}
% ----------------------------------------------------------------

Table~\ref{tab:conc-claim-validation} lists the chapter's main claims,
the evidence behind each, and the boundary within which it holds.  The
purpose is to keep each conclusion within the evidence reported in the
thesis.

{\small
\begin{longtable}{@{}p{0.18\textwidth}p{0.27\textwidth}p{0.25\textwidth}p{0.20\textwidth}@{}}
\caption{Validation mapping between thesis claims, supporting
evidence, and scope boundaries}
\label{tab:conc-claim-validation} \\
\toprule
\textbf{Claim area} & \textbf{Supporting evidence} &
\textbf{Validated conclusion} & \textbf{Scope boundary} \\
\midrule
\endfirsthead
\caption*{Table~\thetable{} (continued): Validation mapping between
thesis claims, supporting evidence, and scope boundaries} \\
\toprule
\textbf{Claim area} & \textbf{Supporting evidence} &
\textbf{Validated conclusion} & \textbf{Scope boundary} \\
\midrule
\endhead
\bottomrule
\endfoot
Correctness &
Selected comparisons with Trino, Oracle, and hand-written
\texttt{pandas} reported in Chapter~\ref{chap:evaluation}. &
The reported reference queries produce the same rows and measure
values across the compared systems. &
Selected supported scenarios only. \\
\midrule
Feature support &
The coverage tables and reference tests in
Chapter~\ref{chap:evaluation}, supported by the design in
Chapters~\ref{chap:methodology} and~\ref{chap:implementation}. &
The system supports a practical row-pattern subset suitable for
DataFrame-centered analysis. &
This is not full SQL:2016 or vendor coverage. \\
\midrule
Performance &
The repeated 30-cell benchmark to 2.22\,M rows and the separate
scalability sweep to 227.9\,M rows. &
The engine has the lowest time on the evaluated cells and shows
near-linear measured scaling. &
One host, selected patterns, and different repetition protocols for
the main and scalability studies. \\
\midrule
Memory use &
Incremental-peak and total-footprint measurements in both benchmarks. &
The measurements show the memory cost of the evaluated workloads. &
Empirical measurements, not a formal upper bound. \\
\midrule
Developer effort &
Lines of detection code in the three worked examples. &
The SQL form uses fewer lines and removes manual traversal logic. &
Code comparison only, not a user study. \\
\midrule
Architecture &
The five-stage design and its automated tests. &
The implementation separates parsing, planning, compilation, and
execution concerns. &
The system remains a \texttt{pandas}-centered prototype. \\
\end{longtable}
}

Three validity limits remain.  \emph{Internal validity} depends on the
reference outputs and their normalization.  \emph{Construct validity}
depends on row equality, time, throughput, memory, and line count
representing the claims made from them.  \emph{External validity} is
limited by one workstation, two performance datasets, and the selected
patterns.
The main benchmark uses twenty measured runs per cell; the larger
sweep uses five.  Its recorded engine and retained database results
were not collected in one simultaneous rerun.  Trino and
Oracle were not expert-tuned, so other configurations may narrow the
measured gaps.

% ----------------------------------------------------------------
\section{Limitations}
\label{sec:conc-limitations}
% ----------------------------------------------------------------

The following limitations define how far the results can be
generalized.

\textbf{Supported subset, not exhaustive SQL.}
The engine targets a substantial but bounded slice of
\texttt{MATCH\_RECOGNIZE}: \texttt{PARTITION BY},
\texttt{ORDER BY}, \texttt{PATTERN}, \texttt{DEFINE}, selected
\texttt{MEASURES}, supported skip and output modes, navigation,
subsets, exclusion blocks, anchors, and \texttt{PERMUTE}.  It does not
claim full SQL:2016 coverage or parity with Trino or Oracle.
General ordering expressions, vendor-specific extensions, and other
unsupported expression forms remain outside the validated scope.

\textbf{Complex pattern and quantifier interactions.}
The automata design handles ordinary sequences, alternation, grouping,
bounded repetition, and common quantifiers.  Deeply nested expressions
with several greedy or unbounded quantifiers can blow up the state and
transition count, and \texttt{PERMUTE} is especially costly because its
semantics are factorial.  The engine therefore uses explicit NFA,
DFA, and \texttt{PERMUTE} guards.  If complete construction exceeds the
state budget, compilation fails.  It does not execute or cache an incomplete
automaton.  This preserves correctness, but a very large pattern may be
rejected.

Predicates set a second limit.  A \texttt{DEFINE} that cannot become
column operations may need the exact search, which is the engine's
slowest path.  The evaluated running \texttt{AVG} form is updated
incrementally, but this does not cover every state-dependent
expression.  Exact search also has an explicit search budget.  Its
exhaustion is reported as an error rather than a changed result
(Section~\ref{subsec:eval-worst-cases}).

\textbf{Measure and aggregate scope.}
The measure evaluator supports common functions---\texttt{SUM},
\texttt{COUNT}, \texttt{MIN}, \texttt{MAX}, \texttt{AVG}, statistical
and collection/string aggregates, conditional aggregates, and supported
\texttt{RUNNING}/\texttt{FINAL} forms.  Arbitrary user-defined
aggregates and every possible expression combination remain out of
scope---a natural limit of an in-memory engine with its own
expression evaluator rather than a full SQL engine.

\textbf{Performance boundaries.}
The core benchmark is predictable from 100\,K to 2.22\,M rows, and a
separate sweep confirms near-linear execution time to 227.9\,M rows.
Measured engine memory also increased close to linearly
on one workstation and one CPU core.  Wider schemas, many small
partitions, and highly ambiguous patterns were not tested at that
scale.  For large row-local scans, available RAM is the practical
limit.  Pattern compilation and exact search also have explicit
algorithmic guards.

\textbf{Memory behavior.}
Memory stayed within the configured limits, and query memory increased
with input size.  Small deltas still vary with allocator reuse, garbage
collection, result materialization, and the 50\,ms sampling interval.
The reported values are therefore empirical measurements, not a formal
upper bound.  The input DataFrame and result must remain resident, so
the current engine is not an out-of-core system.

\textbf{Optimizer integration.}
The engine runs independently of database optimizers.  It does no
cost-based plan selection, predicate pushdown into external storage,
join reordering, or distributed scheduling, so some plan-level
optimizations that databases offer are simply absent here.

% ----------------------------------------------------------------
\section{Conclusion}
\label{sec:conclusion}
% ----------------------------------------------------------------

Pandas has no native row-pattern primitive: detecting a multi-row
pattern means shifted columns, manual loops, and custom state
handling, or a round trip through an external SQL engine.  This thesis
addressed that gap with a SQL-in-\texttt{pandas} implementation that
runs a supported subset of \texttt{MATCH\_RECOGNIZE} over
DataFrames.

The work makes three contributions.  First, it provides an in-process
Pandas engine with verified semantics for the evaluated R010-style
subset.  Second, it provides an optimized and guarded execution design
for conditions, matching, measures, caching, output construction, and
memory use.  Third, it provides
theoretical and empirical evaluation of correctness, coverage,
usability, performance, and scalability.

The evidence supports three conclusions.  The selected correctness
scenarios agree with Trino, Oracle, and hand-written \texttt{pandas}.
The declarative examples use less detection code than their procedural
baselines.  The main and scalability benchmarks show predictable
growth under the stated protocols, with an increasing in-memory cost.

The approach is a bridge between SQL row-pattern
recognition and DataFrame analytics.  It does not replace mature SQL
engines; it brings a useful slice of their row-pattern power into
Python workflows where the data already lives in \texttt{pandas}.
The conclusion holds within the limits stated in
Section~\ref{sec:conc-limitations}.

% ----------------------------------------------------------------
\section{Future Work}
\label{sec:conc-future}
% ----------------------------------------------------------------

Five directions would extend the work.

\textbf{Broader SQL and dialect coverage.}
The most direct step is wider SQL:2016 and vendor coverage: more
navigation combinations, broader expressions, richer aggregates, fuller
\texttt{RUNNING}/\texttt{FINAL} semantics, and a larger conformance
suite drawn from Trino, Oracle, and the standard.

\textbf{Parallel and distributed execution.}
Independent partitions are a natural unit of parallel work.  The
current sequential partition loop could use independent workers with
separate matcher state while preserving deterministic output order.
Distributed execution on Dask or Spark would also need a rule for a
partition split across chunks.

\textbf{Backend abstraction and columnar acceleration.}
The path is \texttt{pandas}-centered today.  A backend abstraction
could target Polars or Apache Arrow, pushing more predicate evaluation
into compiled vectorized kernels and reducing Python iteration.  An
out-of-core backend could also reduce the resident-memory limit, but it
would have to preserve partition order and matches across batch
boundaries.

\textbf{Optimization and diagnostics.}
The automata pipeline can gain from stronger state pruning, guard
pushdown, better selectivity estimates, and a faster exact search for
state-dependent forms not covered by compiled plans.  This remains a
performance boundary (Section~\ref{subsec:eval-worst-cases}).  Better
\texttt{EXPLAIN} output, match traces, and per-pattern counters would
make each execution path easier to verify.

\textbf{Usability and integration.}
Clearer error messages, notebook-friendly visualizations, reusable test
fixtures, and integration with existing DataFrame query interfaces
would make the engine easier to adopt in exploratory analysis and
teaching.

% ----------------------------------------------------------------
\section{Chapter Summary}
\label{sec:conc-summary}
% ----------------------------------------------------------------

This chapter drew the findings together in a scoped context.  It read
the evaluation in practical terms: supported semantics, developer
effort, performance, and memory.  It also separated measured results
from their limits.

The conclusion is that a practical subset of SQL row-pattern
recognition can run directly over DataFrames.  This is supported by the
reported correctness and performance evidence, not by a claim of full
SQL parity.  Future work follows the remaining limits: wider
conformance, parallel execution, alternative backends, exact-search
optimization, and better diagnostics.
